% Generated draft: requires mathematical and source review.
\documentclass[11pt]{article}
\usepackage[margin=25mm]{geometry}
\usepackage{fontspec}
\setmainfont{Latin Modern Roman}
\usepackage{amsmath,amssymb,mathtools,unicode-math}
\setmathfont{Latin Modern Math}
\usepackage{enumitem,needspace,xurl,hyperref}
\setlength{\parindent}{0pt}
\setlength{\parskip}{5pt}
\setlength{\emergencystretch}{3em}
\newcommand{\N}{\mathbb{N}}
\newcommand{\Z}{\mathbb{Z}}
\newcommand{\Q}{\mathbb{Q}}
\newcommand{\R}{\mathbb{R}}
\newcommand{\C}{\mathbb{C}}
\newcommand{\F}{\mathbb{F}}
\newcommand{\A}{\mathbb{A}}
\newcommand{\PP}{\mathbb{P}}
\newcommand{\E}{\mathbb{E}}
\newcommand{\eps}{\varepsilon}
\newcommand{\dd}{\,\mathrm{d}}
\begin{document}
\section*{UnsolvedMath ID: 30001811}
\textbf{Source problem code:} catalogue entry
\subsection{Definitions and mathematical statement}
Let \(B\) be a spherical building equipped with its standard angular metric \(d_B\); its apartments are unit spherical Coxeter complexes. A nonempty subset \(C\subseteq B\) is called \(\pi\)-convex if, for every \(x,y\in C\) with \(d_B(x,y)<\pi\), the unique geodesic segment from \(x\) to \(y\) lies in \(C\). No condition on geodesics joining points at distance exactly \(\pi\) is imposed by this definition.

Two possible circumradius conventions are
\[
\operatorname{rad}_{C}(C)=\inf_{c\in C}\sup_{x\in C}d_B(c,x),
\qquad
\operatorname{rad}_{B}(C)=\inf_{b\in B}\sup_{x\in C}d_B(b,x).
\]
The first restricts centers to \(C\); the second allows centers anywhere in \(B\). Neither definition requires the infimum to be attained.

\textbf{Source question, with unresolved conventions.} Are convex subsets of spherical buildings with circumradius strictly greater than \(\pi/2\) necessarily subbuildings? Here ``rigid'' means precisely the asserted subbuilding conclusion, rather than an additional rigidity property.

Under the explicit conventions that convexity means \(\pi\)-convexity and circumradius means \(\operatorname{rad}_{C}\), the proposed universal assertion takes the form
\[
\begin{gathered}
\text{For every spherical building }B\text{ and every nonempty }C\subseteq B,\\
C\text{ is }\pi\text{-convex and }\operatorname{rad}_{C}(C)>\frac{\pi}{2}
\quad\Longrightarrow\quad
C\text{ is a subbuilding of }B.
\end{gathered}
\]
This is a conditional transcription of the intended question, not a resolution of its missing conventions. If ambient centers are intended, replace \(\operatorname{rad}_{C}\) by \(\operatorname{rad}_{B}\). The strict inequality excludes circumradius exactly \(\pi/2\).

A subbuilding is understood informally here as a subset carrying a spherical-building structure compatible with that of \(B\). The supplied record does not specify the permitted rank, type, or precise compatibility requirement for this conclusion. It also does not specify closedness of \(C\), thickness or rank restrictions on \(B\), or whether its use of ``convex'' includes a condition at antipodal pairs. These qualifications require clarification before the displayed assertion can be treated as a fully determined formalization. The supplied target concerns arbitrary convex subsets; it does not restrict them to convex subcomplexes.
\subsection{Short English statement}
The question asks whether every convex subset of a spherical building whose circumradius is strictly greater than \(\pi/2\) must itself be a subbuilding. The radius convention, precise convexity convention, and intended meaning of subbuilding need clarification. The supplied record labels the problem partially solved and reports that the unrestricted convex-subset question remains open; these are supplied dated assessments requiring independent review.
\subsection{Sources}
\begin{itemize}
\item Local UnsolvedMath problem registry, permanent ID \(30001811\), source problem code \texttt{catalogue entry}, title ``Rigidity of Large-Radius Convex Subsets in Buildings,'' in the collection ``Oberwolfach Reports Open Problems.'' The supplied proposed year is 2011; no individual proposer is supplied.
\item \emph{Arithmetic Groups vs. Mapping Class Groups: Similarities, Analogies and Differences} (2011), Oberwolfach Reports, pp.\ 1637--1708, DOI 10.4171/owr/2011/30. Supplied original-report link: \url{https://doi.org/10.4171/owr/2011/30}. The background describes Oberwolfach Reports as documenting mathematical workshops and their open-problem discussions. No problem-specific page is supplied.
\item Supplied literature source labeled ``Ramos--Cuevas convex subcomplexes'': \url{https://arxiv.org/abs/1408.3017}. The record supplies no further bibliographic details for this reference.
\item Literature assessment dated exactly \texttt{2026-08-21}: supplied dated curation, not independently verified here. It labels the problem ``partially solved'' and states that the Kleiner--Leeb dichotomy is proved for convex subcomplexes in important building types, including type A and the relevant type D cases; that the center conjecture is settled for thick buildings; and that the unrestricted convex-subset question remains open. The background describes this as literature triage derived from the supplied Oberwolfach curation and requests independent verification of time-sensitive status and cited claims.
\item The registry supplies the statement-status label \texttt{corrected\_verified} and reports that the mathematical content was checked against the official Oberwolfach report, with notation, context, and question form normalized without changing the source-level claim. This is a supplied verification claim; the cited sources have not been independently inspected for this draft.
\end{itemize}
\end{document}
